% Workaround para criação de capítulo não-numerado alinhado à margem
\chapter*{}
\noindent
\phantomsection{\MakeUppercase{\textbf{Introdução}}}
\addcontentsline{toc}{chapter}{INTRODUÇÃO}

% \newline
% \newline

%\section{Motivação}

\vspace{1.4cm}


Introdução ao trabalho, usando citações sempre.
    
\vspace{0.5cm}
\textbf{Objetivo}
\vspace{0.5cm}

Falar sobre o objetivo do trabalho.

Falar sobre a relevância e importância do trabalho - venda seu peixe. 

\vspace{0.5cm}
\textbf{Estrutura do Trabalho}
\vspace{0.5cm}

O Capítulo 1 apresenta uma revisão bibliográfica.

O Capítulo 2 explora a teoria e os conceitos envolvidos.

O Capítulo 3 demonstra a metodologia utilizada no projeto e o que foi desenvolvido.

O Capítulo 4 descreve os resultados obtidos.

O Capítulo 5 expõe a conclusão do trabalho.




\vspace{10mm} %5mm vertical space

